% Personalization Strategies Subsection
\subsection{Personalization Strategies}

Personalization strategy represents a fundamental design choice in wearable seizure monitoring systems. Global models train on population data and apply to all patients. Patient-specific models train on individual patient data. Mixed approaches combine population patterns with individual adaptation.

\subsubsection{Global Models}

Five studies use global model approaches. \textcite{spahrDeepLearningbasedDetection2025} trained a global model on 384 patients with tunable parameters. \textcite{fineDetectionKeyAutomated2025} trained on 15 patients and tested on 3 patients. \textcite{dongTwoLayerEnsembleMethod2022} trained on 68 patients with 10-fold cross-validation. \textcite{elemamAutomatedValidatedTool2025} trained on 198 questionnaire responses and 30 HRV recordings. \textcite{meiselMachineLearningWristband2020} trained a global model tested with LOSO validation.

Global models offer advantages for deployment. They enable on-device processing without requiring individual patient training. They can be validated with larger datasets. However, they may not capture individual seizure patterns.

\textcite{meiselMachineLearningWristband2020} used LOSO validation to assess global model generalizability. The study achieved 62\% patient success (43 of 69 patients above chance). This result indicates that global models do not work equally well for all patients.

\subsubsection{Patient-Specific Models}

Four studies use patient-specific approaches. \textcite{stirlingForecastingSeizureLikelihood2021} trained individual models with weekly retraining for 11 patients. \textcite{nasseriAmbulatorySeizureForecasting2021} trained separate models for each of 6 patients. \textcite{odeDevelopmentEpilepticSeizure2023} trained anomaly detection models at the 99\% confidence level for each of 66 patients. \textcite{singhrathoreDevelopmentMultimodalMachine2024} presented a single-patient case study.

Patient-specific models achieve higher individual success rates. \textcite{stirlingForecastingSeizureLikelihood2021} achieved 100\% patient success for hourly forecasting. \textcite{nasseriAmbulatorySeizureForecasting2021} achieved 83\% patient success (5 of 6 patients). These rates are substantially higher than the 62\% achieved by global models with LOSO validation.

Patient-specific approaches require sufficient individual data. \textcite{stirlingForecastingSeizureLikelihood2021} required a mean of 14.6 months of data per patient. \textcite{nasseriAmbulatorySeizureForecasting2021} required 60+ days per patient. This data requirement creates a cold start problem where new patients must wait weeks or months before the system becomes useful.

\subsubsection{Mixed Approaches}

\textcite{vielufSeizureMonitoringCombined2025} used a mixed approach. The study combined population-level 24-hour harmonic patterns with individual diary data. This hybrid achieved 82\% sensitivity and 67\% specificity. The approach may balance the advantages of global and patient-specific methods.

\subsubsection{Subject-Split Validation}

\textcite{reintjesECGBasedDetectionEpileptic2025} used subject-split validation. The study trained on one set of patients and tested on different patients. This approach assesses generalization to new patients. The results showed poor generalization with FPR ranging from 1.91/h to 39.75/h depending on method.

\textcite{nasseriAmbulatorySeizureForecasting2021} used temporal split validation. Early data from each patient served as training. Later data served as testing. This approach assesses temporal stability rather than generalization to new patients.

\subsubsection{Comparison by Strategy}

Global models achieve competitive performance in detection applications. \textcite{spahrDeepLearningbasedDetection2025} achieved 96\% sensitivity with 0.005/h FPR using a global approach. This represents the best FPR among all detection studies.

Patient-specific models achieve higher individual success in forecasting applications. \textcite{stirlingForecastingSeizureLikelihood2021} achieved 100\% patient success compared to 62\% for the global model tested by LOSO in \textcite{meiselMachineLearningWristband2020}.

\textbf{Trade-off identified.} Patient-specific models achieve higher individual success but require weeks to months of individual data. Global models enable immediate deployment but may not work well for all patients. Mixed approaches may offer a middle path.

\subsubsection{Computational Considerations}

Global models enable on-device deployment. \textcite{spahrDeepLearningbasedDetection2025} achieved 112 ms inference time suitable for real-time processing. On-device processing enables continuous monitoring without cloud connectivity.

Patient-specific models often require cloud processing. \textcite{stirlingForecastingSeizureLikelihood2021} used smartphone-based processing with weekly retraining. \textcite{nasseriAmbulatorySeizureForecasting2021} used cloud-based processing. Cloud-based approaches may experience connectivity issues and latency.

\textbf{Key limitation.} Detection studies rarely report patient-level success. Only two of eight detection studies provide individual patient results. This reporting gap prevents assessment of which patients benefit from global detection models.
